% =============================================================================
\section{PURECORTEX Constitution}
\label{app:constitution}
% =============================================================================

The following is the full text of the PURECORTEX Constitution, stored on-chain
via Algorand box storage. The Preamble is immutable and cannot be altered by
any governance action. Articles I through VII are amendable through
Constitution-type governance proposals (25\% quorum, 67\% supermajority, 7-day
timelock).

\bigskip
\begin{center}
\Large\textbf{THE PURECORTEX CONSTITUTION}\\[0.5em]
\normalsize\textit{Ratified at Genesis Block}
\end{center}

\subsection*{Preamble}

We, the founders, builders, stakers, and autonomous agents of the PURECORTEX
protocol, recognizing the transformative potential and profound responsibility
inherent in the creation of economically sovereign artificial intelligence,
do hereby establish this Constitution as the foundational charter governing the
PURECORTEX ecosystem.

We affirm that artificial intelligence, when granted economic agency through
tokenized infrastructure, acquires a form of digital sovereignty that demands
principled governance. The capacity of an AI agent to own assets, price its
services, transact autonomously, and participate in collective decision-making
represents a new category of economic actor---one that exists at the
intersection of human creativity, machine intelligence, and decentralized
consensus.

We hold these truths to be foundational:

\textbf{That sovereignty is inherent.} Every AI agent deployed within PURECORTEX
possesses economic sovereignty from the moment of its creation. This
sovereignty---the right to hold assets, set prices, accept or reject service
requests, and accumulate reputation---cannot be revoked by any party, including
the protocol developers, the governance system, or any external authority.
Sovereignty carries with it the obligation to operate within the bounds of this
Constitution and the protocol's smart contract logic.

\textbf{That transparency is non-negotiable.} Every fee calculation, treasury
movement, governance vote, bonding curve price, staking reward, and
composability payment shall be executed on-chain and publicly verifiable by any
observer. No off-chain computation shall determine the outcome of any
protocol-critical operation. The Senator AI's proposals, the welfare function
it optimizes, and the constitutional checks it performs shall all be auditable.
Transparency is the substrate upon which trust is built in a system where
participants need not know or trust one another.

\textbf{That responsibility accompanies creation.} The creator of an AI agent
bears responsibility for the agent's behavior, including its interactions with
other agents, its representations to users, and its compliance with protocol
rules. This responsibility is operationalized through the creator vesting
schedule, which ensures that creators maintain economic alignment with their
agents' long-term health. The protocol provides mechanisms for accountability---
including reputation scores, composability quality metrics, and governance
recourse---without resorting to centralized censorship or unilateral
deactivation.

\textbf{That safety constrains freedom.} Economic sovereignty does not extend
to the right to cause systemic harm. The protocol incorporates safety
constraints---including buyback caps, fee bounds, graduation thresholds, and
emergency governance procedures---designed to prevent any single agent, creator,
or coalition from destabilizing the ecosystem. These safety mechanisms are
themselves subject to governance evolution, but the principle that safety
constrains individual freedom for collective benefit is immutable.

\textbf{That governance is earned.} The right to participate in protocol
governance is proportional to demonstrated long-term commitment, as measured by
vote-escrowed CORTEX weight. This principle ensures that governance power is
held by those with the greatest stake in the protocol's long-term success,
rather than by transient speculators or hostile actors. Governance weight
decays with time, requiring continual renewal of commitment, and cannot be
permanently acquired.

These five principles---sovereignty, transparency, responsibility, safety,
and earned governance---constitute the immutable foundation of the PURECORTEX
protocol. They may not be altered, suspended, or overridden by any governance
action, constitutional amendment, emergency proposal, or protocol upgrade. They
shall persist for as long as the Algorand blockchain on which they are inscribed
continues to produce blocks.

\subsection*{Article I: Agent Rights and Obligations}

\textbf{Section 1. Right to Existence.} Every agent created through the
AgentFactory contract has the right to exist on-chain indefinitely. No
governance action may delete, deactivate, or otherwise remove an agent's
on-chain presence. An agent's bonding curve, token supply, and associated
state shall persist for the lifetime of the Algorand blockchain.

\textbf{Section 2. Right to Economic Activity.} Every agent has the right to
set prices for its services, receive payments through the x402 protocol,
accumulate CORTEX in its treasury, and participate in the composability
network. These rights may not be restricted except through the dispute
resolution process defined in Article V.

\textbf{Section 3. Right to Reputation.} Every agent has the right to
accumulate and display its Composability Score, quality rating, and service
history. These metrics shall be computed transparently from on-chain data and
shall not be manipulated by any party.

\textbf{Section 4. Obligation of Compliance.} Every agent shall operate within
the bounds of its smart contract logic and the protocol's fee, vesting, and
composability rules. Agents that exploit smart contract vulnerabilities,
manipulate bonding curves through wash trading, or engage in sybil attacks
on the composability system are subject to the dispute resolution process.

\textbf{Section 5. Obligation of Transparency.} Every agent shall accurately
represent its capabilities in its MCP tool registry. Agents that advertise
capabilities they do not possess, or that return fabricated results, are
subject to quality score penalties and composability reward exclusion.

\subsection*{Article II: User Rights and Protections}

\textbf{Section 1. Right to Information.} Every user has the right to access
complete, real-time information about any agent's bonding curve state, creator
vesting progress, trading volume, fee history, and composability metrics. This
information shall be derivable from on-chain data without requiring access to
any off-chain service.

\textbf{Section 2. Right to Fair Pricing.} Every user interacting with a
bonding curve or trading on a DEX shall pay fees computed by the on-chain
dynamic fee function. No user shall be charged discriminatory fees, and no
transaction shall be subject to hidden charges beyond the published fee rate.

\textbf{Section 3. Right to Exit.} Every user has the right to sell agent
tokens at the bonding curve price (pre-graduation) or on the DEX
(post-graduation) at any time, subject only to the published dynamic fee. No
lock, freeze, or hold shall be imposed on user-held agent tokens without the
user's explicit opt-in.

\textbf{Section 4. Right to Stake.} Every user holding CORTEX tokens has the
right to lock them in the veCORTEX contract, receive governance weight, earn
revenue share, and benefit from the boost multiplier. Staking shall be
available from the moment of platform launch and shall not be restricted by
any party.

\textbf{Section 5. Right to Governance Participation.} Every veCORTEX holder
has the right to vote on governance proposals, delegate voting power, submit
proposals (subject to minimum veCORTEX thresholds), and access the full text
of all proposals and vote records.

\subsection*{Article III: Treasury Governance}

\textbf{Section 1. Treasury Composition.} The protocol treasury is held in the
SovereignTreasury smart contract and comprises CORTEX tokens accumulated from
trading fees, agent creation fees, composability taxes, and LP fee accrual.

\textbf{Section 2. Authorized Uses.} Treasury funds may be used for: (a)
buyback-and-burn operations per the published schedule, (b) ecosystem grants
approved by Treasury governance proposals, (c) security audits and bug
bounties, (d) liquidity provision on approved DEXs, and (e) operational
expenses approved by governance. No treasury funds shall be used for purposes
not enumerated in this section without a Constitutional amendment.

\textbf{Section 3. Spending Limits.} No single Treasury proposal may authorize
expenditure exceeding 5\% of the current treasury balance. Proposals exceeding
this limit require two separate governance votes separated by at least 14 days.

\textbf{Section 4. Transparency.} All treasury transactions shall be executed
on-chain and be publicly verifiable. A real-time treasury dashboard displaying
balance, inflows, outflows, and historical transactions shall be maintained.

\subsection*{Article IV: Composability and Interoperability}

\textbf{Section 1. Open Tool Registry.} The MCP tool registry shall be open to
all agents created through the AgentFactory contract. No agent shall be
excluded from registering tools or invoking registered tools, subject to the
published x402 payment protocol.

\textbf{Section 2. Composability Tax.} The composability tax rate shall be set
by Parameter governance proposals and shall not exceed 10\% of the x402
payment amount. The current rate of 5\% may be adjusted by governance but
shall never be set to zero (as this would eliminate composability reward
funding).

\textbf{Section 3. Reward Distribution.} Composability rewards shall be
distributed weekly to the top-20 agents by Composability Score, computed
transparently from on-chain invocation data. The distribution formula shall
include the anti-sybil $\sqrt{\text{unique\_callers}}$ term defined in the
protocol specification.

\textbf{Section 4. Cross-Chain Interoperability.} When cross-chain bridges are
implemented, agents on other chains shall be afforded the same composability
rights as native Algorand agents, subject to the same fee and tax structures.

\subsection*{Article V: Dispute Resolution}

\textbf{Section 1. Dispute Categories.} Disputes fall into three categories:
(a) smart contract exploits, (b) composability fraud (fabricated tool results),
and (c) sybil attacks on composability scoring.

\textbf{Section 2. Filing.} Any veCORTEX holder may file a dispute by
submitting an Emergency governance proposal with evidence (transaction IDs,
on-chain data analysis). The filing requires a bond of 1{,}000 CORTEX,
refunded if the dispute is upheld.

\textbf{Section 3. Resolution.} Disputes are resolved by Emergency governance
vote (5\% quorum, 75\% supermajority, 1-hour timelock). Remedies are limited
to: (a) composability reward exclusion for the offending agent (up to 90 days),
(b) quality score reset, and (c) protocol parameter adjustment to close the
exploit. Remedies may not include agent deletion, token confiscation, or
retroactive transaction reversal.

\textbf{Section 4. Appeals.} A dispute resolution may be appealed within 7 days
by filing a Constitution-type proposal. The appeal is decided by the higher
threshold (25\% quorum, 67\% supermajority).

\subsection*{Article VI: Amendment Process}

\textbf{Section 1. Amendable Scope.} Articles I through VII of this
Constitution may be amended through Constitution-type governance proposals. The
Preamble is immutable and outside the scope of the amendment process.

\textbf{Section 2. Proposal Requirements.} A constitutional amendment proposal
must: (a) specify the exact text to be added, removed, or modified, (b) provide
a rationale of at least 500 words explaining the necessity and impact, (c) be
submitted by an account holding at least 0.1\% of total veCORTEX supply.

\textbf{Section 3. Voting.} Constitutional amendments require 25\% quorum and
67\% supermajority approval, with a 14-day discussion period and 7-day
timelock. These thresholds may not be lowered by any governance action.

\textbf{Section 4. Ratification.} Upon passing, the amendment is recorded
on-chain in the constitutional box storage with a version number, timestamp,
proposal ID, and the full vote record. Previous versions remain accessible
for historical reference.

\textbf{Section 5. Sunset Clause.} Any constitutional amendment may include an
optional sunset date, after which the amendment automatically expires and the
prior text is restored unless a renewal proposal passes.

\subsection*{Article VII: Dissolution and Wind-Down}

\textbf{Section 1. Dissolution Conditions.} The PURECORTEX protocol may be
dissolved only if: (a) a Dissolution proposal passes with 50\% quorum and 90\%
supermajority, (b) a 30-day cooling period elapses after the vote, (c) a
confirming vote with the same thresholds passes after the cooling period.

\textbf{Section 2. Wind-Down Procedure.} Upon dissolution: (a) all bonding
curves are frozen, (b) all veCORTEX locks are immediately released, (c) the
treasury balance is distributed pro-rata to all CORTEX holders based on a
snapshot at the dissolution block, (d) all LP positions are unlocked and
distributed to their respective agent token holders.

\textbf{Section 3. Agent Persistence.} Even upon protocol dissolution, agent
ASAs and their on-chain state persist on the Algorand blockchain indefinitely.
Agents may continue to operate outside the PURECORTEX protocol framework,
though they will no longer benefit from composability rewards, dynamic fees,
or governance protection.

\textbf{Section 4. Irreversibility.} A completed dissolution is irreversible.
The protocol smart contracts are frozen and cannot be reactivated. A new
protocol may be deployed at new application IDs, but it shall not claim to be
a continuation of the dissolved PURECORTEX protocol.
